\section{Non-linear effects: variance contribution predicts the ranking}
\label{sec:nonlinear}

Data: $y = a + b^2 + c^3$, with $a = 1 + U_a$, $b = 1 + U_b$, $c = 1 + U_c$ and
$U_a, U_b, U_c$ independent $U(0,1)$ draws, $n = 10000$. (The draws are
\emph{separate} per feature: a single shared $U$ would make the three features
collinear and the experiment vacuous.)
Each term's variance contribution grows with its degree, so
Section~\ref{sec:background-strands} predicts
$\mathrm{vip}(a) < \mathrm{vip}(b) < \mathrm{vip}(c)$; the forest confirms it,
$1751 < 8166 < 39632$ (a $1:4.7:22.6$ ratio). A linear model still returns non-zero
coefficients, but they are uninterpretable: the residuals are visibly non-normal.
